\begin{statementbox}
\begin{center}\begin{tikzpicture}[font=\small,>=stealth,line join=round]
\draw[gray!40,rounded corners=3pt] (0,0) rectangle (14.2,8.3);\node at (7.1,7.700000000000001) {左右点不必同高，但连线的中点始终是中心};
\begin{scope}[shift={(1.491304347826087,2.4166666666666665)},x=1.3043478260869565cm,y=0.5111111111111111cm]\draw[->,gray] (-0.3,0)--(4.3,0) node[right]{$x$};\draw[->,gray] (0,-1.5)--(0,7.5) node[above]{$y$};\draw[blue,thick,solid,domain=0.4:3.6,samples=180] plot(\x,{(\x-2)^3+3});\fill[orange!85!black] (1,2) circle(2.2pt);\fill[orange!85!black] (3,4) circle(2.2pt);\draw[orange!85!black] (1,2)--(3,4);\draw[gray,dashed] (2,-1.2)--(2,7.2);\draw[gray,dashed] (-.2,3)--(4.2,3);\fill[violet] (2,3) circle(2.6pt);\end{scope}\node[anchor=west,align=left,text width=5.0cm] at (8.4,4.0) {$f=(x-2)^3+3$\\[9pt]$P(3,4),\ Q(1,2)$\\[9pt]$\frac{3+1}{2}=2,\quad\frac{4+2}{2}=3$\\[9pt]中点 $C(2,3)$\\[9pt]$f(2+t)+f(2-t)=6$};
\node[text width=12.6cm,align=center] at (7.1,.65) {这条具体曲线示范一般公式；严格判定仍要求所有点都有对应点。};
\end{tikzpicture}\end{center}


\textbf{一句话记住：中心对称，横坐标取平均，纵坐标也取平均。}
两个平均数合起来，就是对称中心的坐标。

\textbf{\textcolor{CStatement}{这个结论解决什么问题？}}

不必把整条函数曲线画出来，也不必求出每个点，只用一条等式就能判断图像
是不是\textbf{绕某个点转半圈后仍与原图重合}；反过来，已知对称中心，还能快速求函数值。

\medskip
\textbf{\textcolor{CStatement}{先看条件（两个都不能少）}}
设 $y=f(x)$ 的定义域为 $D$，准备判断图像是否关于点 $C(a,b)$ 中心对称。
\begin{enumerate}[leftmargin=*,itemsep=0.3em]
  \item \textbf{横坐标要能成对出现：}对每个实数 $t$，$a+t\in D$ 当且仅当 $a-t\in D$。
  \item \textbf{成对函数值的和固定：}当 $a+t,a-t\in D$ 时，必须总有
  \[
     \boxed{\ f(a+t)+f(a-t)=2b\ }.
  \]
\end{enumerate}
\textbf{结论：}以上两条同时成立，\textbf{当且仅当} $y=f(x)$ 关于 $(a,b)$ 中心对称。
“当且仅当”就是\textbf{从图形能推出公式，从公式也能保证图形}。

\medskip
\end{statementbox}\clearpage\begin{statementbox}\textbf{\textcolor{CStatement}{考试中更顺手的写法}}
\[
  \boxed{\ f(x)+f(2a-x)=2b\quad(x\in D)\ },
\]
这里仍须确保 $x\in D\Longrightarrow 2a-x\in D$。
两个输入 $x$ 与 $2a-x$ 的平均数是 $a$，两个输出的平均数是 $b$。

\medskip
\textbf{\textcolor{CStatement}{与奇函数的联系}}
令 $g(t)=f(a+t)-b$，则 $g(-t)=-g(t)$；也就是说，\textbf{把中心 $(a,b)$ 移到原点后，就变成奇函数}。

\medskip
\textbf{\textcolor{CStatement}{一个容易被忽略的细节}}
若 $a\in D$，代入 $t=0$ 得 $2f(a)=2b$，所以 $f(a)=b$。
但若 $a\notin D$，对称中心仍可能存在，\textbf{不一定在函数图像上}。
\end{statementbox}
\clearpage
\begin{statementbox}
\textbf{同一个判定的常见等价写法：两个平均数读出中心}
\begin{center}\begin{tikzpicture}[font=\small,>=stealth,line join=round]
\draw[gray!40,rounded corners=3pt] (0,0) rectangle (14.2,8.3);\node at (7.1,7.700000000000001) {先看一对点：中心就是连线的中点};
\begin{scope}[shift={(1.3,1.5)},x=1.4cm,y=1.1cm]\draw[->,gray] (0,0)--(4.7,0) node[right]{$x$};\draw[->,gray] (0,0)--(0,4.8) node[above]{$y$};\draw[blue,thick] (1,2)--(3,4);\fill[orange!85!black] (1,2) circle(2.4pt);\fill[orange!85!black] (3,4) circle(2.4pt);\fill[violet] (2,3) circle(2.5pt);\node[below right=7pt] at (1,2){$Q(1,2)$};\node[right=7pt] at (3,4){$P(3,4)$};\node[below right=8pt,violet] at (2,3){$C(2,3)$};\end{scope}\node[anchor=west,align=left] at (8.4,4.5){横坐标平均为2\\[14pt]纵坐标平均为3\\[14pt]$\overrightarrow{CP}=-\overrightarrow{CQ}$};
\node[text width=12.6cm,align=center] at (7.1,.65) {不指定整条函数曲线；一对点的对称，不能证明整张图像对称。};
\end{tikzpicture}\end{center}


以下先取定义域为 $\mathbb R$，$a,b,p,q,c$ 为固定实数，每行等式均对所有实数自变量成立：
\begin{center}
\renewcommand{\arraystretch}{1.45}
\begin{tabular}{p{.56\linewidth}|p{.28\linewidth}}
等式的写法 & 一个对称中心\\\hline
$f(a+t)+f(a-t)=2b$ & $(a,b)$\\
$f(x)+f(2a-x)=2b$ & $(a,b)$\\
$f(2a-x)=2b-f(x)$ & $(a,b)$\\
$f(p+t)=c-f(q-t)$ & $\left(\dfrac{p+q}{2},\dfrac c2\right)$
\end{tabular}
\end{center}
\textbf{只需一个读式规则：}最后一行移项后是 $f(p+t)+f(q-t)=c$。两个输入的和恒为 $p+q$，两个输出的和恒为 $c$，\textbf{分别除以 $2$}就得到中心坐标。令 $u=p+t$，便是 $f(u)+f(p+q-u)=c$，与原判定完全相同。

图中一对点的横坐标为 $3,1$，纵坐标为 $4,2$，两个平均数为 $2,3$，所以这对点的中心是 $(2,3)$。
\textbf{其他定义域：}须保证每个 $u\in D$ 都有 $p+q-u\in D$，且相应函数值之和对全部配对恒为 $c$；不能用图中一对点代替全域验证。
\end{statementbox}
